Для сравнения моделей хранения необходимо зафиксировать предметную область.
Модель должна быть близкой к реальным системам видеонаблюдения, но
без лишних деталей, которые бы влияли на хранение и анализ метаданных.

\subsection{Общая характеристика системы видеонаблюдения}

Современные системы видеонаблюдения используются не только для визуального
контроля территории. Часто они служат источником различных данных, которые могут
быть использованы для анализа событий.
В общем случае такая система включает камеры, зоны наблюдения,
серверы обработки или хранения, программное обеспечение для управления
потоками данных и анализа видеозаписей. Согласно руководствам по
проектированию CCTV-систем, при построении системы важно учитывать
качество изображения, размещение камер, а также механизмы передачи,
обработки и хранения данных \cite{dhs_cctv_handbook}.


Видеопоток остается основным типом данных, которые собирает система
видеонаблюдения. Однако помимо видеоданных системы генерируют множество
метаданных и событий, которые могут быть не менее важными для анализа.
Например, умные камеры могут распознавать объекты, определять их тип (человек, автомобиль, животное),
отслеживать их перемещение, фиксировать время и место событий.
Эти данные используют для поиска закономерностей, обнаружения аномалий и
других задач, которые выходят за рамки простого просмотра видеозаписей.

В рамках данной работы система видеонаблюдения будет рассматриваться не как
полноценный программно-аппаратный комплекс охраны, а как ее часть. 
Поэтому исследование фокусируется на метаданных и
исключает хранение самого видеопотока.

Система строится вокруг хранения и анализа событий,
полученных от большого количества камер. Камеры фиксируют происходящее в
зонах наблюдения, а встроенная в камеру система аналитики формирует события:
обнаружение движения, появление человека или автомобиля, пересечение виртуальной линии,
нахождение объекта в запрещенной зоне, потерю сигнала с камеры и другие ситуации.
Для каждого события сохраняются метаданные, необходимые для последующего поиска,
фильтрации и анализа.

Эти метаданные загружаются в СУБД, которая обеспечивает их хранение, индексацию и
доступ для аналитических задач.

\subsection{Сценарий работы системы}

Рассмотрим сценарий: имеется охраняемая территория, разделенная на
несколько зон наблюдения. В каждой зоне установлены камеры. Каждая камера
передает данные в локальное хранилище, расположенное на охраняемой территории.
Данные внутри этого хранилища обрабатываются и анализируются, итоговая информация
сохраняется на удаленном сервере в базе данных.
Система не сохраняет на удаленном сервере видеозаписи, но
фиксирует сведения о состоянии камер и событиях,
полученных в результате внутренней аналитики.

Основной поток данных образуют события видеонаблюдения. Событие представляет
собой факт, зафиксированный камерой или аналитической системой
(внутри камеры или на локальной системе), который может быть так или иначе
сохранен в локальном хранилище охраняемой территории.
Оттуда, в любом случае, итоговое событие передается на общий сервер, где
сохраняется в базе данных.

Например, система может обнаружить движение, потерю сигнала, появление объекта
или пересечение заданной линии в кадре. При этом для каждого события
сохраняются дата и время возникновения, зона, одна или несколько камер,
зафиксировавших событие, тип события, уровень значимости, уверенность
аналитического алгоритма и дополнительные параметры.

Нас интересуют именно эти метаданные, приходящие потоком на удаленный сервер.
Важно зафиксировать их структуру, объем, частоту и другие характеристики,
чтобы использовать равные условия для сравнения моделей данных.

\image{scenario.png}{Сценарий работы системы}{1}
